\scanpage{5}\mspage{\handwritten{4}}

\deleted{pleinement fidèle (resp.} une équivalence). Même énoncé si on met dans le coup un \deleted{Anneau \unsure{loc.}} système local d'anneaux sur $X'$. Ehfin, $f$ induit \deleted{une} \deleted{isomorphisme}\inserted{\deleted{foncteur pleinement fidèle (resp.} une équivalence)} entre la catégorie des coefficients locaux sur $C$ et celle des coéff. locaux sur $C'$ \deleted{(?? peut-être \inserted{cette propriété doit} \unsure{cela} doit-il être mis dans les hypothèses préliminaires sur $f$ ?).} Principe de démonstration: on commence par prouver ce dernier résultat, en notant que \deleted{dans le ca} si un topos est localement hom. trivial, il est loca connexe et loc. simplement connexe, d'où une bonne théorie du $\handwritten{\pi}_{\unsure{o|0}}$ et du $\handwritten{\pi}_1$ (qui sont ici discrets), et on est ramené \deleted{pour} à un cas particulier du critère d'homotopisme de Artin et Mazur, savoir un critère \inserted{cohomologique} pour qu'un homomorhisme de groupes $G\handwritten{\to}H$ (ici des groupes $\handwritten{\pi}_1$ de $X,X'$) soit un isomorphisme: il doit induire des isomorphismes sur les $H^{\unsure{o|0}}$ et $H^1$ (y inclus dans le cas non commutatif...). On prouve la pleine fidélité en se ramenant par la suite spectrale aux cas des $\mathrm{Ext}^i$ de coéfficients locaux. Par la suite spectrale encore, cela résultera du fait suivant: si $X^{\deleted{'}}$ est \unsure{homalement|honalement|localement} hom. trivial, alors la catégorie des faisceaux abéliens loc. constants est stable par $\underline{\mathrm{Ext}}^i$, et le foncteur $M\handwritten{\mapsto}M_X$ des \deleted{faisceaux constants} $(\mathrm{Ab})$ dans $X_{ab}$ commute aux dits $\mathrm{Ext}^i$.

En fait, $X$ et $X'$ étant loc.hom\deleted{p}triviaux, les conditions suivantes sur $f$ seront équivalentes:

a) $f$ est un homotopisme, i.e. induit pour tout système local (pas néc commutatif) sur $X'$ un isomorphisme sur les $H^i$\deleted{pour \unsure{$i<n$|$i\le n$}}.

b) \deleted{Même chose ave} $f$ induit une équivalence $D^b_{lc}(X')\handwritten{\longrightarrow}D^b_{lc}(X)$.

a\handwritten{'}) $f$ induit une équivalence entre la catégorie des systèmes locaux abéliens sur $X'$ et $X$, et des isomorphismes sur les $H^i$ correspondants (\unsure{donc} on ne prend ici que des coéfficients commutatifs).

\margin{\handwritten{?}} J'ignore si on peut dans a) se borner aux systèmes locaux commutatifs \deleted{i.e. dans}.

L'équivalence entre a) et b) fournit une première justification ou motivation pour définir des types d'homotopie via la catégorie $D\handwritten{^{b}_{lc}(X)}$,

% archive sheet number: 4
